\clearpage\begingroup\SourceLeaf{57b}\HistoricalHead{}{L}{PHYSIK. BRIEFWECHSEL.}
\begin{center}
[4.]\\[.7em]
\textsc{Gauss} an \textsc{Encke}. Göttingen, 15. August 1842.
\end{center}
\noindent .\ .\ .\ .\ .\ .\ .\ .\ .\ .\ .\ .\ .\ .\ .\ .\ .\ .\ .\ .\ .
\par\noindent\hspace*{\parindent}Für Ihre Mittheilung der Artillerie Versuche bin ich Ihnen sehr dankbar.
\textsc{Bessels} Aburtheilung ist unstreitig zu schroff. Es giebt ohne Zweifel viele
Fälle, wo man Beobachtungszahlen, auch ohne sie mit einer Theorie bemeistert
zu haben, mit Nutzen einer Interpolation unterwerfen kann, in so fern
man von der wirklichen Zuverlässigkeit aller jener Beobachtungszahlen eine
völlige Ueberzeugung hat. Von der anderen Seite ist nicht zu leugnen, dass
gerade dieser Ueberzeugung zumahl bei etwas verwickeltern Gegenständen
die volle Lebendigkeit fehlen kann, wenn man nicht ihren Zusammenhang unter
sich mit einer Einsicht in ihre Theorie durchdringt. Ich bekenne, dass die
mir mitgetheilten Zahlen mir doch vielen Zweifel an der Zuverlässigkeit, trotz
der Uebereinstimmung unter sich, zurücklassen, namentlich scheint es mir gar
nicht denkbar, dass die Schussweite von \(15^0\) bis \(20^0\) Elevation mehr zunehmen
sollte als von \(10^0\) bis \(15^0\), wie bei 5 Pfund Ladung die mir mitgetheilten
Zahlen ergeben: 1303,1 1729,9 2212,0.

Um sich darüber ein Urtheil zu bilden, müsste man alle Details der Versuche
genau kennen. Uebrigens gestehe ich, dass mir scheint, aus Versuchen,
bie \textls[120]{dloss} auf diese Art angestellt werden, könne nicht viel herauskommen.
\par\noindent\begin{minipage}[t]{.42\textwidth}\vspace{0pt}\centering
\begin{tikzpicture}[x=.008cm,y=-.008cm,line width=.55pt]
\draw (99,382) .. controls (215,237) and (327,85) .. (546,43);
\draw (159,330) -- (211,246);
\draw (250,184) -- (291,139);
\draw (50,426) -- (90,376) -- (105,388) -- (65,439) -- cycle;
\draw (33,451) -- (50,430) -- (59,438) -- (42,459) -- cycle;
\draw (27,465) -- (35,454) -- (41,459) -- (32,470);
\draw (24,465) -- (33,461) -- (37,470) -- (30,476) -- cycle;
\end{tikzpicture}
\end{minipage}\hfill
\begin{minipage}[t]{.55\textwidth}\vspace{0pt}\setlength{\parindent}{1.2em}
\noindent Man sollte, meine ich, auf Mittel denken,
jedem Schuss \textls[120]{mehrere} Resultate abzugewinnen,
als bloss die Schussweite. Ich
dächte z. B., dass es wohl thunlich wäre,
mehrere Tangenten der Trajectorie mit
grosser Schärfe zu erhalten, wenn hinterwerts
in schicklicher Entfernung Beobachter
mit Höhenmessungsinstrumenten aufgestellt
wären. Vielleicht liessen sich sogar von
der Seite her brauchbare Beobachtungen
\end{minipage}
\clearpage\SourceLeaf{58}\HistoricalHead{53}{R}{GAUSS AN ENCKE.}
\noindent machen, wenn die Aufstellung der Beobachter und Instrumente schicklich
arrangirt wäre. \textsc{Schumacher} schrieb mir vor 2 Jahren[\textsuperscript{*})]\begingroup\renewcommand{\thefootnote}{}\footnotetext{[\textsuperscript{*}) Siehe \textit{Briefwechsel zwischen Gauss und Schumacher} III, Altona 1861, S. 402.]}\endgroup, dass in der Sammlung
des russischen Generalstabs in St. Petersburg 155 \textsc{Ertel}sche Theodolithen
sind. Ich weiss nicht, ob der preussische eben so viel besitzt. Aber
mit einem oder ein \textls[120]{Paar Dutzend} Theodolithen würde ich jeden Schuss
verfolgen lassen, und ich dächte, dass wenn man es recht anfinge, die Trajectorie
in vielen Punkten sich müsse ergreifen lassen. Ich würde nur
stärkere Elevationen gebrauchen. Ist man einmal Herr der Curven und
der Bewegungsgesetze, so kann man das nöthige für kleine Elevationen besser
durch Rechnung als durch Versuch finden.\quad .\ .\ .\ .\ .\ .\ .\ .\ .\ .\ .\ .\ .\ .\ .\ .\ .\ .\ .\ .
\clearpage\endgroup
